\SourceLeaf{315}%
Nachdem man dahin gelangt ist, dass bei dem Gebrauch von fünfziffrigen Logarithmen der Werth von \(T\) sich nicht mehr ändert, berechnet man \(\lambda\) nach der Formel
\[
\lambda=r\sin T\sec B
\]
und führt dann eine neue Rechnung mit sieben Decimalen, wobei man die logarithmischen Correctionen vermittelst der Formeln
\[
\begin{aligned}
\log b&=\log\beta+\mu\lambda\lambda+\tfrac12\mu\tau\tau\\
\log t&=\log\tau+\mu rr+\tfrac12\mu\tau\tau
\end{aligned}
\]
zuzieht und \(B=B^0-\frac12b\), \(T=T^0-\frac12t\) setzt. Eine nochmalige Wiederholung wird in der Regel dieselben, oder kaum merklich geänderte Resultate wiedergeben, und dann erst wird auch noch die Berechnung von \(l\) nach der Formel
\[
\log l=\log\lambda-\tfrac12\mu rr+\tfrac12\mu\lambda\lambda
\]
beigefügt. Um die Schnelligkeit der Annäherung (die hauptsächlich von der Kleinheit von \(r\) abhängt), an einem Beispiele zu zeigen, setze ich die Hauptmomente der Rechnung für den Übergang von dem Dreieckspunkte Brocken zu dem Punkte Inselsberg hieher. Es ist diess die grösste Dreiecksseite in dem Hannoverschen Dreieckssystem, viel grösser, als sonst bei trigonometrischen Operationen vorzukommen pflegen.

Bei der nach den Grundlagen der ersten Abhandlung bearbeiteten conformen Darstellung auf der Kugelfläche ist die Breite des Brockens \(B^0=51^0\,46'\,3''\,6345\); das Azimuth der Seite Brocken-Inselsberg \(T^0=5^0\,42'\,21''\,7704\); der Logarithm dieser Seite in Toisen \(=4{,}7353929\) oder in Theilen des Halbmessers \(=8{,}22018543\), oder in Bogensecunden, wie bei unsern Formeln vorausgesetzt ist, \(\log r=3{,}5346106\). Setzt man zuerst \(T=5^0\,42'\), so wird
\[
\begin{aligned}
\beta&=3408''\\
B&=51^0\,17'\,40''\\
\tau&=424''
\end{aligned}
\]
und folglich ein genäherterer Werth
\[
T=5^0\,38'\,50''
\]
